\chapter[Джанака]{Радість самопізнання (Джанака)}

[2.1] Справді, я чистий і сповнений спокою. Я --- Свідомість по той бік природної причинності. Як довго я дозволяв омані мучити мене!

[2.2] Як я освітлюю це тіло, так само освітлюю й Всесвіт. Отже, увесь Всесвіт --- мій; або, навпаки, ніщо мені не належить.

[2.3] Тепер, коли я перестав ототожнювати себе з тілом і речами, завдяки великій Благодаті відкривається Істина.

[2.4] Як хвилі, піна й бульбашки не відрізняються від води, так усе, що постає з Я, не відрізняється від нього.

[2.5] Уважно розглянь тканину --- й побачиш лише нитки. Уважно поглянь на творіння --- й побачиш лише Я.

[2.6] Як цукор із цукрової тростини пронизаний солодкістю, так усе це, вияв Я, цілковито пронизане мною.

[2.7] Через незнання себе виникає ілюзія. Завдяки самопізнанню вона зникає. Мотузка --- не змія, хоч може нею здаватися.

[2.8] Світло --- моя істинна природа, не більше й не менше. Всесвіт проявляється в моєму спогляданні.

[2.9] Ілюзія невігластва виникає в мені так само, як змія постає з мотузки, марево води --- на сонячному обрії, а срібло --- у перламутрі.
